% SLICEKW1 research note
\documentclass[11pt]{article}
\usepackage[margin=.9in]{geometry}
\usepackage{amsmath,amssymb,amsthm}
\newtheorem{proposition}{Proposition}
\begin{document}

\begin{center}
{\Large SLICEKW1: Communication structure of the exact-layer 3-UNSAT slice}\\[4pt]
\textbf{Date: 2026-10-04}\\
\textbf{Status: OPEN / CANDIDATE}\\
\textbf{P versus NP status: OPEN}
\end{center}

\section{Setup}

Retain the SLICEPAD1 notation.  Let \(S_{q,t}\) be the monotone slice extension of clause-presence
\(\mathrm{UNSAT}_{3,q}\), and let
\[
\mathcal U_{q,t}
=
\left\{
A\in\binom{\mathcal C_q}{t}:
A\text{ is UNSAT}
\right\}.
\]
The negative exact-layer inputs are
\[
\mathcal S_{q,t}
=
\left\{
B\in\binom{\mathcal C_q}{t}:
B\text{ is SAT}
\right\}.
\]

The monotone Karchmer--Wigderson relation restricted to the exact layer is
\[
R_{q,t}(A,B)
=
\{C\in\mathcal C_q:C\in A\setminus B\},
\qquad
A\in\mathcal U_{q,t},\ B\in\mathcal S_{q,t}.
\]
Since \(S_{q,t}(A)=1\) and \(S_{q,t}(B)=0\), the set \(A\setminus B\) is nonempty.

\section{A simple low-communication protocol}

\begin{proposition}
There is a deterministic protocol for \(R_{q,t}\) using \(q+O(1)\) communicated bits, apart from the standard convention for naming the final output coordinate.
\end{proposition}

\begin{proof}
Bob has a satisfiable formula \(B\).  He selects any satisfying assignment
\(\beta\in\{0,1\}^q\) and sends \(\beta\) to Alice using \(q\) bits.
Since \(A\) is unsatisfiable, some clause \(C\in A\) is falsified by \(\beta\).
Because \(\beta\) satisfies every clause of \(B\), that clause \(C\) cannot belong to \(B\).
Hence \(C\in A\setminus B\), as required.
\end{proof}

This protocol is an important negative control.  Ordinary deterministic communication complexity
of the exact-layer monotone KW relation is at most \(O(q)\).  Therefore an ordinary communication-depth lower bound cannot by itself yield the desired superpolynomial circuit-size lower bound.

The auxiliary search relation
\[
(A,\beta)\longmapsto
C\in A\cap B_\beta
\]
is useful for constructing protocols, where \(B_\beta\) is the block of clauses falsified by \(\beta\).
It should not be identified with the full KW relation without proving a reduction between the two games.

\section{Why DAG-like protocols are the correct next model}

For formulas, tree-like KW protocols capture depth/formula phenomena.  For general monotone
circuits with sharing, the relevant exact characterization is via DAG-like communication protocols.
The modern formulation of Garg--G\"o\"os--Kamath--Sokolov characterizes monotone circuit size by the least size of a DAG-like protocol solving the monotone KW relation.

Thus the next target is not
\[
\mathrm{CC}(R_{q,t})=\omega(q),
\]
which is impossible by the protocol above, but rather a size lower bound
\[
\mathrm{DAGKWsize}(R_{q,t})=M_q^{\omega(1)}.
\]

The sparse slice \(t=6q\) is preferred initially because it retains NP/coNP hardness while positive
inputs contain only \(O(q)\) clauses.  The exact assignment-block geometry remains
\[
|B_\alpha|=\binom q3,\qquad
|B_\alpha\cap B_\beta|
=
\binom{q-d_H(\alpha,\beta)}3.
\]

\section{Strength of the desired conclusion}

The SLICEPAD1 padding theorem shows that exact-\(t\) SAT is NP-complete for both \(t=6q\) and the middle slice.  Therefore the complement of the total slice function \(S_{q,t}\) is an NP-complete language: below the layer it accepts automatically, on the layer it is SAT, and above the layer it rejects.

Valiant's slice theorem converts a superpolynomial monotone lower bound for \(S_{q,t}\) into a
superpolynomial unrestricted circuit lower bound for \(S_{q,t}\).  Complementing the output gate
changes unrestricted circuit size by at most a constant.  Consequently the desired theorem would imply
\[
\mathrm{NP}\not\subseteq\mathrm{P/poly},
\]
which is strictly stronger than merely \(P\ne NP\).

This is a useful calibration.  The slice route is genuine, but it is not a shortcut around general
nonuniform circuit lower bounds.

\section{Hostile implications}

\begin{itemize}
\item[] The \(q\)-bit satisfying-assignment protocol blocks any plan based only on large ordinary KW communication complexity.
\item[] The no-constant-weight-monotone-embedding theorem from SLICEPAD1 blocks importing a full-cube lower bound by a monotone projection.
\item[] Any successful argument must control DAG sharing, not only tree-like protocol depth.
\item[] Any lifting theorem must respect the exact-weight domain.  A reduction from a full Boolean cube to one fixed layer cannot itself be coordinatewise monotone unless constant.
\end{itemize}

\section{Next concrete target}

Construct a structured subrelation
\[
R^\star_{q}\subseteq R_{q,6q}
\]
whose positive formulas are exact-layer UNSAT transversals and whose negative formulas admit controlled
satisfying assignments, and then prove one of:

\begin{enumerate}
\item a superpolynomial lower bound on DAG-like protocol size for \(R^\star_q\); or
\item a lifting theorem from a query lower bound to DAG-like protocol size that preserves the exact-weight domain.
\end{enumerate}

The Hamming association-scheme identity
\[
|B_\alpha\cap B_\beta|
=
\binom{q-d_H(\alpha,\beta)}3
\]
is the principal available symmetry.  The first falsification test should determine whether this symmetry also permits a compact DAG protocol, in which case the route should be retired or further restricted.

\section{Final status}

SLICEKW1 narrows the communication target but proves no lower bound.  Ordinary communication complexity is explicitly easy.  The surviving circuit-size target is a DAG-like protocol size lower bound.  \(P\ne NP\) remains \textbf{OPEN}.

\begin{thebibliography}{4}
\bibitem{Valiant1986}
L. G. Valiant, Negation is Powerless for Boolean Slice Functions, \emph{SIAM Journal on Computing} 15(2):531--535, 1986.

\bibitem{Wegener1985}
I. Wegener, On the Complexity of Slice Functions, \emph{Theoretical Computer Science} 38:55--68, 1985.

\bibitem{GGKS2020}
A. Garg, M. G\"o\"os, P. Kamath, and D. Sokolov,
Monotone Circuit Lower Bounds from Resolution,
\emph{Theory of Computing} 16(1), 2020.

\bibitem{Lokam2003}
S. V. Lokam,
Graph Complexity and Slice Functions,
\emph{Theory of Computing Systems} 36:1--88, 2003.
\end{thebibliography}
\end{document}
